\documentclass[10pt,a4paper]{amsart}
\usepackage[margin=19mm]{geometry}
\usepackage{amsmath,amssymb,mathtools}
\usepackage[scr=rsfs,cal=euler]{mathalfa}
\usepackage{xfrac}
\usepackage[unicode,hidelinks,bookmarks=false]{hyperref}
\newcommand{\ie}{i.e.\ }
\newcommand{\cf}{cf.\ }
\newcommand{\resp}{respectively\ }
\newcommand{\Subsection}{\subsection}
\providecommand{\nobreakdash}{\nobreak\hbox{-}}
\setlength{\emergencystretch}{2em}
\multlinegap0pt
\usepackage{mathtools}
\usepackage{amssymb}
\usepackage{xcolor}
\usepackage{cancel}
\usepackage[shortlabels]{enumitem}
\newtheorem{theorem}{Theorem}
\newtheorem{lemma}{Lemma}
\newcommand{\Emin}{E_\textup{min}}
\newcommand{\NE}{\mathcal{NE}(\mathbb{R})}
\newcommand{\R}{\mathbb{R}}
\newcommand{\boC}{\mathcal{C}}
\newcommand{\boE}{\mathcal{E}}
\newcommand{\boX}{\mathcal{X}}
\newcommand{\bu}{\mathbf{u}}
\renewcommand{\d}{\,\mathrm{d}}
\title{Momentum--mass normalized dark--bright solitons to one dimensional Gross--Pitaevskii systems}
\author{Salvador L\'opez-Mart\'inez}
\date{Source: arXiv:2508.19216v3 (CC BY 4.0)}
\begin{document}
\maketitle

\noindent\emph{Source and licence.} Momentum--mass normalized dark--bright solitons to one dimensional Gross--Pitaevskii systems. Version arXiv:2508.19216v3 (CC BY 4.0), distributed by arXiv under the Creative Commons Attribution 4.0 International licence, http://creativecommons.org/licenses/by/4.0/, as recorded in the arXiv licence metadata for this exact version and shown on the abstract page https://arxiv.org/abs/2508.19216v3. Full licence notice: \texttt{licenses/arxiv-2508.19216-CC-BY-4.0.txt}.\par\smallskip

\noindent\textit{Editorial scope.} Counted unit: the article's theorem thm:properties (source0.tex lines 337--358). The article's complete original proof of it is delivered with it (source0.tex lines 593--656, one \texttt{proof} environment of 64 lines, which the article places away from the statement). No mathematics is written, repaired or re-derived: the statement and the proof are the article's own lines, in the article's own notation, and no retained line is substituted or re-flowed.  Retained uncounted as the source-local context the proof needs: lines 313--315; lines 520--522; lines 552--565.  Nothing was omitted from any retained span, so the package carries no omission record.  No retained line contains a citation, so the wrapper carries no bibliography and no bibliography entry.  No retained line contains a figure environment, an image-inclusion command or drawing code, so no image is delivered and the package declares no asset dependency.  Editorial formatting only, in addition: an AMS \texttt{amsart} preamble that restates the article's own declarations and macros used by the retained lines, namely the environments \texttt{theorem}, \texttt{lemma} on separate counters, together with the standard packages those lines need.  The retained environments print in this document as Theorem 1, Lemma 1 in order of appearance, whereas the article prints the same items with the numbering of its own full document.  The complete native source of the article as distributed in the arXiv e-print for this exact version is bundled unchanged as \texttt{sources/183942/source0.tex}.
\begin{equation}\label{eq:explicitsol}
    \bu_q(x)=\sqrt{\frac{2-c^2}{{2}}}\tanh\left(\frac{\sqrt{2-c^2}}{2}x\right)-\frac{c}{\sqrt{2}}i.
\end{equation}

\begin{lemma}\label{lemma:sameinf}
    For every $\alpha>0$, $\beta\geq 0$, $q>0$, $m\geq 0$, the following identity holds: \[\Emin(q,m)=\inf\{E(u,v):\, (u,v)\in \boX^0_{q,m}\}.\]
\end{lemma}

\begin{lemma}\label{lemma:smallE}
    Let $\alpha>0$, $\beta\geq 0$, $q\in(0,\pi/2)$, and $m\geq 0$. Then, 
    \begin{equation}\label{eq:eminineq}
        \Emin(q,m)\leq \Emin(q,0)<\sqrt{2}q.
    \end{equation}
    Assuming in addition 
    \begin{equation}\label{eq:alphabeta}
        m\beta< 2\alpha\int_\R(1-|\mathbf{u}_q|^2)\d x, \quad m>0,
    \end{equation}
    where $\bu_q$ is defined in \eqref{eq:explicitsol}, it follows that
    \begin{equation}\label{eq:keyineq}
        \Emin(q,m)<\Emin(q,0).
    \end{equation}
\end{lemma}

\begin{theorem}\label{thm:properties}
	Let $\alpha>0$ and $\beta > 0$ satisfy $\alpha^2\leq\beta$. 
	\begin{enumerate}
		\item \label{item:monotoneq} For every $m\geq 0$, the function $q\mapsto \Emin(q,m)$ is nondecreasing in $[0,\infty)$ and $\Emin(0,m)=0$. In particular, $\Emin(q,m)\geq 0$ for all $q\in [0,\pi/2)$, $m\geq 0$.
		\item \label{item:monotonem} For every $q\geq 0$, the function $m\mapsto \Emin(q,m)+\alpha m/2$ is nondecreasing in $[0,\infty)$. 
		\item \label{item:decreasingm} For $0\leq  q_1\leq q_2<\pi/2$ and $0\leq m_1\leq m_2$, one has
			\begin{equation}\label{lipsineq}
				\Emin(q_2,m_2)\leq \Emin(q_1,m_1)+\sqrt{2}(q_2-q_1).
			\end{equation}
            In particular, for every $q\in [0,\pi/2)$, the function $m\mapsto \Emin(q,m)$ is nonincreasing in $[0,\infty)$, and
            \begin{equation}\label{eq:energeticallyfav}
            \Emin(q,m)<\Emin(q,0),\quad\text{for all }q\in (0,\pi/2),\, m>0.
            \end{equation}
		\item  \label{item:lipschitz} The function $(q,m)\mapsto\Emin(q,m)$ is Lipschitz in $[0,\pi/2)\times[0,\infty)$.
        \item \label{item:subadd} For $q_1\in [0,\pi/2)$, $q_2\in [0,\pi/2)$, $m_1\geq 0$, and $m_2\geq 0$, it follows 
    \begin{equation}\label{eq:subadditive}
    \Emin(q_1+q_2,m_1+m_2)\leq\Emin(q_1,m_1)+\Emin(q_2,m_2),
    \end{equation}
    and the equality in \eqref{eq:subadditive} holds if, and only if, $(q_1+q_2)(q_1+m_1)(q_2+m_2)=0$.
    
	\end{enumerate}
\end{theorem}

\begin{proof}[Proof of Theorem~\ref{thm:properties}~\ref{item:monotoneq}-\ref{item:lipschitz}]
	To prove \ref{item:monotoneq}, let $m, q_1,$ and $q_2$ be nonnegative real numbers such that $q_1<q_2$, and let $\varepsilon>0$. Take $(u,v)\in\boX_{q_2,m}$ be such that $E(u,v)<\Emin(q_2,m)+\varepsilon$. Since $u\in\NE$, there is a lifting $u=\rho e^{i\theta}$.  Let $\lambda=q_1/q_2\in [0,1)$, and $u_\lambda=\rho e^{i\lambda\theta}$. It is clear that $p(u_\lambda)=q_1$, and
\[\Emin(q_1,m)\leq E(u_\lambda,v)\leq E(u,v)\leq\Emin(q_2,m)+\varepsilon.\]
Letting $\varepsilon$ tend to zero, we deduce that the function $q\mapsto \Emin(q,m)$ is nondecreasing in $[0,\infty)$. 

{On the other hand, given $(u,v)\in\boE(\R)\times H^1(\R)$, Young inequality and $\alpha^2\leq\beta$ imply
\begin{equation*}
\frac14(1-|u|^2)^2 + \frac{\beta}{4}v^4 -\frac{\alpha}{2}(1-|u|^2) v^2 \geq \frac14\Big(1-\frac{\alpha^2}{\beta}\Big)(1-|u|^2)^2 \geq 0.
\end{equation*}
Hence, $E(u,v)\geq 0$, so that $\Emin(q,m)\geq 0$ for all $q\in [0,\pi/2)$, $m\geq 0$. 

In addition, let $m\geq 0$ and $v\in H^1(\R)$ such that $\|v\|_2^2=m$. For every $\lambda>0$, define $v_\lambda(x)=\sqrt{\lambda}v(\lambda x)$, for all $x\in\R$. It is clear that $\|v_\lambda\|_2^2=m$ for all $\lambda>0$, and
\[E(1,v_\lambda)=\int_\R\left(\frac{\lambda^2}{2}(v')^2+\frac{\beta\lambda}{4}v^4\right)\d x.\]
Therefore $E(1,v_\lambda)\to 0$ as $\lambda\to 0$. In particular, since $(1,v_\lambda)\in\boX_{0,m}$ for all $\lambda>0$, it follows that $\Emin(0,m)=0$.} This finishes the proof of \ref{item:monotoneq}. 

Let us now prove \ref{item:monotonem}. To do so,  let $q, m_1,$ and $m_2$ be nonnegative real numbers such that $m_1<m_2$, and let $\varepsilon>0$. Take $(u,v)\in\boX_{q,m_2}$ satisfying that $E(u,v)<\Emin(q,m_2)+\varepsilon$. Let $\lambda=m_1/m_2\in [0,1)$, and $v_\lambda=\sqrt{\lambda}v$. It is clear that $\|v_\lambda\|^2_2=m_1$, and
\begin{align*}
	\Emin(q,m_1)&+\frac{\alpha m_1}{2}\leq E(u,v_\lambda)+\frac{\alpha m_1}{2}
\\ 
&=E(u,0)+\int_\R\left(\frac{\lambda}{2}(v')^2+\frac{\beta\lambda^2}{4}v^4-\frac{\alpha}{2}v^2(1-\lambda|u|^2)\right)\d x +\frac{\alpha m_2}{2}
\\
&\leq E(u,v)+\frac{\alpha m_2}{2}<\Emin(q,m_2)+\frac{\alpha m_2}{2}+\varepsilon.
\end{align*}
We conclude again by letting $\varepsilon$ tend to zero.


In order to show \eqref{lipsineq}, let $\varepsilon>0$, $0\leq q_1\leq q_2<\pi/2$, $0\leq m_1\leq m_2$, and let us distinguish three different cases:

\emph{Case 1:} Let us first assume that $0<q_1<q_2$. By Lemma~\ref{lemma:sameinf}, there exists $(u_1,v_1)\in\boX_{q_1,m_1}^0$ such that $E(u_1,v_1)<\Emin(q_1,m_1)+\varepsilon$. Analogously, denoting $s=q_2-q_1\in (0,\pi/2)$ and $t=m_2-m_1\geq 0$, there exists $(u,v)\in\boX_{s,t}^0$ such that $E(u,v)<\Emin(s,t)+\varepsilon$. Moreover, Lemma~\ref{lemma:smallE}  implies that $E(u,v)<\sqrt{2}(q_2-q_1)+\varepsilon$. Let us write $u_1=\rho_1 e^{i\theta_1}$ and $u=\rho e^{i\theta}$. There is no loss of generality (by applying a translation) in assuming that $1-\rho_1$ and $1-\rho$ have disjoint supports, and the same holds for $\theta'_1$ and $\theta'$, and for $v_1$ and $v$. Hence, defining 
\[\rho_2=\rho_1+\rho-1,\quad\theta_2=\theta_1+\theta,\quad u_2=\rho_2 e^{i\theta_2},\quad v_2=v_1+ v,\]
we deduce that $(u_2,v_2)\in \boX_{q_2,m_2}$ and
\[\Emin(q_2,m_2)\leq E(u_2,v_2)=E(u_1,v_1)+E(u,v)<\Emin(q_1,m_1)+\sqrt{2}(q_2-q_1)+2\varepsilon.\]
Since $\varepsilon$ is arbitrary, \eqref{lipsineq} follows in this case. 

\emph{Case 2:} Assume now that $q_1=q_2$. Since \eqref{lipsineq} is trivially satisfied for $q_1=q_2=0$, we may assume in addition that $q_1=q_2>0$. Arguing as for the previous case, let $(u_1,v_1)\in\boX_{q_1,m_1}^0$ be such that $E(u_1,v_1)<\Emin(q_1,m_1)+\varepsilon$. In addition, let $v\in\boC_0^\infty(\R)$ be such that $\|v\|_2^2=m_2-m_1$ and $E(1,v)<\varepsilon$ (one may find such a function $v$ by rescaling a given function with mass $m_2-m_1$, as in the proof of item \ref{item:monotoneq} or the proof of Lemma~\ref{lemma:smallE}). Applying a translation if necessary, assume that $v$ and $v_1$ have disjoint supports, and define $v_2=v_1+v$. Then, $\|v_2\|_2^2=m_2$ and
\begin{align*}
\Emin(q_2,m_2)&=\Emin(q_1,m_2)\leq E(u_1,v_2)
\\
&=E(u_1,v_1)+\int_\R\left(\frac12(v')^2+\frac{\beta}{4}v^4-\frac{\alpha}{2}v^2(1-|u_1|^2)\right)\d x
\\
&\leq E(u_1,v_1)+E(1,v)<\Emin(q_1,m_1)+2\varepsilon.    
\end{align*}
Thus, \eqref{lipsineq} follows again.

\emph{Case 3:} If $q_1=0<q_2$, {Lemma~\ref{lemma:smallE}} implies
\[\Emin(q_2,m_2)<\sqrt{2}q_2=\Emin(q_1,m_1)+\sqrt{2}(q_2-q_1).\]
Hence, \eqref{lipsineq} follows one more time.

The fact that $\Emin$ is nonincreasing in $m$ is nothing but Case 2 above. Moreover, fixing $q\in(0,\pi/2)$ and $m>0$, it follows from Lemma~\ref{lemma:smallE} the existence of $m_0\in (0,m)$ such that
\[\Emin(q,m)\leq\Emin(q,m_0)<\Emin(q,0).\]
This concludes the proof of \ref{item:decreasingm}.

We finally prove \ref{item:lipschitz}. Let $q_1,q_2\in [0,\pi/2)$ and $m_1,m_2\in [0,\infty)$. Assume that $q_1<q_2$ and $m_2<m_1$, being the remaining possibilities dealt with analogously. By \ref{item:monotoneq}, \ref{item:monotonem}, and \eqref{lipsineq}, we derive
\begin{align*}
&|\Emin(q_2,m_2) -\Emin(q_1,m_1)|\leq \frac{\alpha}{2}(m_1-m_2)+|\Emin(q_2,m_1)-\Emin(q_1,m_1)|
\\
&+\Big|\Emin(q_2,m_2)+\frac{\alpha m_2}{2}-\Emin(q_2,m_1)-\frac{\alpha m_1}{2}\Big|
\\
&= \Emin(q_2,m_1)-\Emin(q_1,m_1)+\Emin(q_2,m_1)-\Emin(q_2,m_2)+\alpha(m_1-m_2)
\\
&\leq\sqrt{2}(q_2-q_1)+\alpha(m_1-m_2).
\end{align*}
This shows that $\Emin$ is a Lipschitz function. 
\end{proof}

\end{document}
